% =====================================================================
% HPC_DEMO — Non-Convex Polygon Packing
% Class-Scale Parallel Optimization Study + N=200 Hero Run
% Chapter-Based Technical Report
% =====================================================================

\documentclass[11pt]{report}

% -------------------- Packages --------------------
\usepackage[margin=1in]{geometry}
\usepackage{amsmath, amssymb, amsthm}
\usepackage{booktabs}
\usepackage{graphicx}
\usepackage{subcaption}
\usepackage{float}
\usepackage{xcolor}
\usepackage{hyperref}
\usepackage{enumitem}
\usepackage{microtype}
\usepackage{mathtools}

\hypersetup{
  colorlinks=true,
  linkcolor=black,
  urlcolor=blue,
  citecolor=black
}

% -------------------- Macros --------------------
\newcommand{\methodA}{\textbf{MS}}
\newcommand{\methodB}{\textbf{ER-MS}}
\newcommand{\methodC}{\textbf{PT}}

\newcommand{\Lbest}{L_{\text{best}}}
\newcommand{\Emin}{E_{\min}}
\newcommand{\EPSFEAS}{\texttt{EPS\_FEAS}}

\newcommand{\figdir}{figures}

% -------------------- Title --------------------
\title{
\textbf{HPC\_DEMO}\\
Class-Scale Parallel Optimization Study\\
+ N=200 Hero Run\\[0.5em]
\large Non-Convex Polygon Packing Under Fixed Resource Constraints
}
\author{Joel Amir Dario Maldonado Tänori}
\date{\today}

\begin{document}

\maketitle
\tableofcontents
\clearpage

% =====================================================================
\chapter{Introduction}

\section{Purpose of the Study}

This work evaluates three parallel optimization strategies for minimizing
the side length $L$ of a square containing $N$ identical non-convex polygons.

This is a \textbf{demonstrative algorithms study} under strict class-project
constraints:

\begin{itemize}
    \item Maximum 20 CPU cores
    \item $\le$ 4 hours wall-clock per method
    \item One N=200 hero run
    \item Clear, interpretable convergence plots
\end{itemize}

\section{Study Structure}

We split the evaluation into:

\begin{itemize}
    \item \textbf{Graph Suite:} $N \in \{5,10,20,50,100\}$
    \item \textbf{Hero Suite:} $N = 200$
\end{itemize}

The hero run is a demonstration of capability, not a statistical comparison.

% =====================================================================
\chapter{Problem Formulation}

\section{Feasibility Definition}

A configuration is feasible if:

\[
\texttt{outside\_penalty} \le \EPSFEAS
\quad \text{and} \quad
\texttt{overlap\_penalty} \le \EPSFEAS
\]

\section{Primary Metric}

\[
\Lbest(t) =
\min_{\tau \le t}
\{ L(\tau) : \text{feasible at time } \tau \}
\]

\section{Secondary Metrics}

\begin{itemize}
    \item Final $\Lbest(T)$
    \item Final bracket width $L_{hi} - L_{lo}$
    \item Success rate
    \item Minimum penalty $\Emin$
\end{itemize}

\section{Hero Efficiency Metric}

\[
\eta = \frac{N A_{poly}}{L^2}
\]

% =====================================================================
\chapter{Experimental Design}

\section{Hardware Constraints}

\begin{itemize}
    \item 20 CPU cores maximum
\end{itemize}

\section{Parallel Layout}

\subsection{Graph Suite}

\begin{itemize}
    \item $R = 10$ threads per run
    \item Two seeds concurrently
\end{itemize}

\subsection{Hero Run}

\begin{itemize}
    \item $R = 20$
    \item Single run
\end{itemize}

\section{Determinism}

\[
\texttt{seed\_r} =
\texttt{splitmix64}(
\texttt{global\_seed} \oplus
\texttt{hash}(N, run\_id, probe\_idx, r))
\]

% =====================================================================
\chapter{Algorithmic Methods}

\section{Method A: Multi-Start (MS)}

Independent simulated annealing workers.

Early stop if any worker becomes feasible.

\section{Method B: Elite Resample Multi-Start (ER-MS)}

Resample frequency:
\[
K_{resample} = 200N
\]

Elite pool:
\[
k' = \lceil 0.25R \rceil
\]

Noise schedule:
\[
\sigma_{pos} = (0.02 - 0.015f)L
\]
\[
\sigma_{\theta} = (5^\circ - 4^\circ f)
\]

\section{Method C: Parallel Tempering (PT)}

Geometric temperature ladder.

Swap frequency:
\[
K_{swap} = 200N
\]

Early stop only if coldest replica feasible.

% =====================================================================
\chapter{Bisection and Scheduling}

\section{Bracket Initialization}

\[
L_{lo} = \sqrt{N A_{poly}}
\]
\[
L_{hi} = \gamma(N) L_{lo}
\]

\section{Graph Suite Schedule}

Base slice times:

\begin{center}
\begin{tabular}{cc}
\toprule
$N$ & $t_{base}$ \\
\midrule
5,10,20 & 6s \\
50 & 10s \\
100 & 15s \\
\bottomrule
\end{tabular}
\end{center}

Stop when:

\[
L_{hi} - L_{lo} \le 10^{-3} L_{hi}
\]

\section{Hero Schedule}

\begin{itemize}
    \item 60s slices (probes 1–6)
    \item 120s slices (7–12)
    \item 240s slices (13+)
\end{itemize}

% =====================================================================
\chapter{Results — Graph Suite}

\section{Convergence Plots}

% \begin{figure}[H]
% \centering
% \includegraphics[width=0.95\linewidth]{\figdir/graphsuite_Lbest_traces.pdf}
% \caption{Convergence traces $\Lbest(t)$ for all seeds and methods.}
% \end{figure}

\section{Bracket Evolution}

% \begin{figure}[H]
% \centering
% \includegraphics[width=0.95\linewidth]{\figdir/graphsuite_bracket_width.pdf}
% \caption{Bracket width over probes.}
% \end{figure}

\section{Success Rates}

Insert summary table here.

\section{Discussion}

Discuss coordination effects, scaling behavior, and feasibility trends.

% =====================================================================
\chapter{Results — N=200 Hero Run}

\section{Run Configuration}

Describe chosen method and pilot behavior.

\section{Convergence}

% \begin{figure}[H]
% \centering
% \includegraphics[width=0.95\linewidth]{\figdir/hero_Lbest_trace.pdf}
% \caption{Hero convergence trace.}
% \end{figure}

\section{Final Metrics}

Report $L$, bracket width, and efficiency $\eta$.

\section{Final Packing}

% \begin{figure}[H]
% \centering
% \includegraphics[width=0.95\linewidth]{\figdir/hero_packing.pdf}
% \caption{Final N=200 packing visualization.}
% \end{figure}

\section{Qualitative Analysis}

Discuss structural characteristics of the packing.

% =====================================================================
\chapter{Limitations}

\begin{itemize}
    \item Different $R$ values between suites
    \item Small seed counts
    \item Hero run is demonstrative only
\end{itemize}

% =====================================================================
\chapter{Reproducibility and Artifacts}

\begin{itemize}
    \item Commit hash
    \item Config files
    \item Seeds
    \item Output CSV files
    \item SVG files
\end{itemize}

\end{document}